In der vorangegangenen Aufgabe wurde die ganze Funktion $f(z)=(e^z-1)/z$
in eine Potenzreihe entwickelt.
\begin{teilaufgaben}
\item Bestimmen Sie die Nullstellen der Funktion $f(z)$.
\item Die reziproke Funktion hat die Potenzreihenentwicklung
\begin{equation}
g(z)
=
\frac{1}{f(z)}
=
\frac{z}{e^z-1}
=
\sum_{n=0}^\infty \frac{B_n}{n!}z^k.
\label{buch:402:greihe}
\end{equation}
Die Koeffizienten $B_n$ heissen die Bernoulli-Zahlen.
Bestimmen Sie die ersten Bernoulli-Zahlen bis $B_4$, indem 
Sie den Ansatz~\eqref{buch:402:greihe}
in die Gleichung
\[
z = (e^z-1) g(z)
\]
einsetzen und die Koeffizienten durch Koeffizientenvergleich
bestimmen.
\item
Bestimmen Sie den Konvergenzradius der Reihe.
\end{teilaufgaben}

\begin{loesung}
\begin{teilaufgaben}
\item
Die Nullstellen von $f(z)$ sind die Nullstellen von $e^z-1$ oder die
Lösungen der Gleichung
\begin{equation}
e^a(\cos b + i \sin b)
=
1.
\label{buch:402:gleichung}
\end{equation}
Daraus liest man ab, dass $\sin b=0$ sein muss,
d.~h. $b\in \pi\mathbb{Z}$.
Für die ungeraden Vielfachen von $\pi$ führen zu einem negativen $\cos b$,
es bleiben daher nur die geraden Vielfachen von $\pi$.
Um die Gleichung~\eqref{buch:402:gleichung} zu erfüllen muss jetzt nur noch
$e^a=1$ sein, d.~h.~$a=0$.
Die Menge der Nullstellen ist daher $2\pi i\mathbb{Z}\setminus\{0\}$.
\item
Wir schreiben Terme bis vierter Ordnung der beiden Reihen und
multiplizieren aus
\begin{align*}
z
&=
\biggl(
1 + \frac{z}{2} + \frac{z^2}{3!} + \frac{z^3}{4!} + \frac{z^4}{5!} + \ldots
\biggr)
\biggl(
B_0 + B_1z + \frac{B_2}{2!}z^2 + \frac{B_3}{3!}z^3 + \frac{B_4}{4!}z^4 + \ldots
\biggr)
\end{align*}
Etwas allgemeiner kann man schreiben
\begin{align*}
z
&=
\biggl(
\sum_{k=0}^\infty \frac{z^k}{(k+1)!}
\biggr)
\biggl(
\sum_{n=0}^\infty \frac{B_n}{n!}z^n
\biggr)
\end{align*}
\item
Die Funktion $g(z)$ hat Pole bei den Nullstellen von $f(z)$.
Die Pole sind $2\pi ik$, $k\in \mathbb{Z}\setminus\{0\}$, daher ist
der nächste Pol zum Nullpunkt der Punkt $z=\pm 2\pi i$ und daher 
ist der Konvergenzradius $=2\pi$.
\qedhere
\end{teilaufgaben}
\end{loesung}

